% Part: methods
% Chapter: proofs
% Section: cant-do-it

\documentclass[../../../include/open-logic-section]{subfiles}

\begin{document}

\olfileid{mth}{prf}{cnt}

\olsection{मला जमत नाही!{}}

कधी ना कधी आपल्या सर्वांना हार मानावीशी वाटते. पण तुम्हाला हे \emph{जमू शकते}. शिक्षक, अध्यापन-सहायक आणि सहाध्यायी तुम्हाला मदत करू शकतात. त्यांच्याकडे मदत मागा!{} अडचणींचा मोठा पेच होऊ नये म्हणून आणि हार मानावीशी वाटल्यास काय करावे यासाठी काही सूचना पुढे दिल्या आहेत.

प्रश्न यशस्वीपणे सोडवता यावेत यासाठी पुढील गोष्टी करा:
\begin{enumerate}
\item \emph{शक्य तितक्या आधी सुरुवात करा.} सत्रभर आपण व्यस्त असतो आणि आपल्यापैकी अनेकांना काम पुढे ढकलण्याची सवय अडवते. गृहपाठाची सुरुवात लवकर करणे ही करता येण्याजोग्या सर्वात उपयुक्त गोष्टींपैकी एक आहे. त्यामुळे अडकलात तरी रडत बसण्यापलीकडे एखादा उपाय शोधण्यासाठी वेळ मिळतो.
\item \emph{सहाध्यायांशी बोला}. तुम्ही एकटे नाही. वर्गातील इतरांनाही अडचणी येत असतील; पण त्यांच्या अडचणी वेगळ्या असू शकतात. सहाध्यायांशी चर्चा केल्यामुळे प्रश्नाकडे पाहण्याची वेगळी दृष्टी मिळू शकते आणि कदाचित मार्ग सापडेल. अर्थात, त्यांचे उत्तर नुसते उतरवू नका: एखादी खूण मागा किंवा कुठे अडकलात हे समजावून पुढची पायरी विचारा. तुम्हाला समजले की त्यांनाही मदत करा. दुसऱ्याला पुढे जाण्यास मदत केल्याने आणि समजावून सांगितल्याने तुमची स्वतःची समजही सुधारते.
\item \emph{मदत मागा.} तुमच्यासाठी अनेक साधने उपलब्ध आहेत. शिक्षक आणि अध्यापन-सहायक तुमच्या मदतीसाठी आहेत आणि तुम्ही यशस्वी \emph{व्हावे अशी त्यांची इच्छा आहे}. प्रश्न सोडवण्यात मदत करणे आणि प्रक्रियेत नेमकी कुठे अडचण येते हे ओळखणे त्यांना शक्य असले पाहिजे.
\item \emph{थोडी विश्रांती घ्या.} अडकलात, तर त्या प्रश्नाकडे खूप वेळ पाहत बसल्यामुळे तसे झाले असण्याची \emph{शक्यता} आहे. थोडी विश्रांती घ्या, चहा प्या किंवा काही वेळ दुसरा प्रश्न सोडवा; मग ताज्या मनाने पहिल्या प्रश्नाकडे परत या. एखादी रात्र झोप घेऊन दुसऱ्या दिवशी पुन्हा पाहा.
\end{enumerate}

या उपायांसाठी सिद्धतेवर काम बराच वेळ आधी सुरू केलेले असणे आवश्यक आहे, हे लक्षात आले का? गृहपाठ देण्याच्या आदल्या दिवशी पहाटे दोन वाजता सिद्धता सुरू केली असेल, तर हे उपाय फार उपयोगी ठरणार नाहीत.

हे निराशाजनक वाटेल; पण सिद्धता करणे हे शेवटी समाधान देणारे आव्हान आहे. काही लोक हेच व्यवसाय म्हणून करतात; म्हणजे त्यात आनंद देणारे काहीतरी असले पाहिजे. जवळजवळ इतर सर्व गोष्टींप्रमाणे प्रश्न सोडवणे आणि सिद्धता करणे यांनाही सराव लागतो. काही सहाध्यायांना हे सोपे वाटत असेल; बहुधा त्यांनी आधीच बराच सराव केला असेल. फार लवकर हार मानू नका.

एखाद्या प्रश्नासाठी वेळ किंवा धीर संपला, तरी ठीक आहे. त्याचा अर्थ तुम्ही मूर्ख आहात किंवा तुम्हाला कधीच जमणार नाही असा नाही. शिक्षक किंवा दुसरा विद्यार्थी यांच्याकडून तो कसा सोडवतात ते समजून घ्या. कुठे चूक झाली किंवा कुठे अडकलात ते ओळखा, म्हणजे पुढे तशी अडचण आली तर ती टाळता येईल. मग उत्तर न पाहता पुन्हा सोडवण्याचा प्रयत्न करा. पुढच्या वेळी आधी सुरुवात करा आणि मदतही लवकर मागा.

\end{document}
